Soit $a$ un réel strictement positif.

On considère la fonction $f$ définie sur l'intervalle $]0~;~+\infty[$ par 

\[f(x) = a \ln (x).\]

On note $\mathcal{C}$ sa courbe représentative dans un repère orthonormé.

Soit $x_{0}$ un réel strictement supérieur à 1.

\medskip

\begin{enumerate}
\item Déterminer l'abscisse du point d'intersection de la courbe $\mathcal{C}$ et de l'axe des abscisses.
\item Vérifier que la fonction $F$ définie par $F(x) =a [x \ln (x) - x]$ est une primitive de la fonction $f$ sur l'intervalle $]0~;~+\infty[$.
\item En déduire l'aire du domaine bleuté en fonction de $a$ et de $x_{0}$.

\end{enumerate}

\begin{center}
\psset{unit=2cm,arrowsize=2pt 3}
\begin{pspicture*}(-1,-1)(4.5,2)
\psaxes[linewidth=1.25pt,Dx=5,Dy=5]{->}(0,0)(-1,-1)(4.5,2)
\pscustom[fillstyle=solid,fillcolor=blue]{\psplot[plotpoints=2000,linewidth=1.25pt,linecolor=blue]{1}{1.8}{x ln 1.3 mul}\psline(1.8,0.764)(1.8,0)(1,0)}
\psplot[plotpoints=2000,linewidth=1.25pt,linecolor=blue]{0.1}{4.5}{x ln 1.3 mul}
\uput[dr](1,0){1}\uput[d](1.8,0){$x_0$}\uput[ul](1.8,0.764){$M$}
\uput[dr](4,1.6){\blue $\mathcal{C}$}\uput[dl](0,0){O}
\uput[d](4.4,0){$x$}\uput[l](0,1.9){$y$}
\end{pspicture*}
\end{center}

On note $T$ la tangente à la courbe $\mathcal{C}$ au point $M$ d'abscisse $x_{0}$.

On appelle $A$ le point d'intersection de la tangente $T$ avec l'axe des ordonnées et $B$ le projeté orthogonal de $M$ sur l'axe des ordonnées.

\begin{center}
\psset{unit=2cm,arrowsize=2pt 3}
\begin{pspicture*}(-1,-1)(4.5,2)
\psaxes[linewidth=1.25pt,Dx=5,Dy=5]{->}(0,0)(-1,-1)(4.5,2)
\psplot[plotpoints=2000,linewidth=1.25pt,linecolor=blue]{0.1}{4.5}{x ln 1.3 mul}
\uput[dr](1,0){1}\uput[d](1.8,0){$x_0$}\uput[ul](1.8,0.764){$M$}
\uput[dr](4,1.6){\blue $\mathcal{C}$}\uput[dl](0,0){O}
\psline[linestyle=dashed](1.8,0.764)(1.8,0)
\uput[d](4.4,0){$x$}\uput[l](0,1.9){$y$}
\psplotTangent{1.8}{4}{x ln 1.3 mul}
\psline(-1,0.764)(4.5,0.764)\uput[ul](0,0.764){B}\uput[ul](0, -0.5){A}
\psframe(0,0.764)(0.1,0.664)\uput[ul](3.1,1.8){$T$}
\end{pspicture*}
\end{center}

\begin{enumerate}
\setcounter{enumi}{3}
\item Démontrer que la longueur AB est égale à une constante (c'est-à-dire à un nombre qui ne dépend pas de $x_{0}$) que l'on déterminera.

\emph{Le candidat prendra soin d'expliciter sa démarche}.
\end{enumerate}
